%% kapa: exam=1 round=25 year=2014 period=1 form=A subject=civil_law subject_ko=민법 seq=32 number=32
%% source: Q-Net 감정평가사 기출문제 (cst00309) artlSeq=5101207 file=extracted/1008/5101207/2014년 제25회 감정평가사 1차시험 기출문제/1교시(A형).hwp
%% conversion: kordoc 4.24.51 -> tools/gen1.py (auto); review=proofread-sonnet; flags=none
\begin{question}{32}
지상권에 관한 설명으로 옳지 않은 것은? (다툼이 있으면 판례에 의함)
\begin{choices}
  \item 지상권자는 지상권을 유보한 채 지상물의 소유권만을 양도할 수 있다.
  \item 당사자가 지료에 관한 합의를 하였더라도 이를 등기하지 않았다면, 당해 지상권자에 대하여 지료지급을 청구할 수 없다.
  \item 구분지상권의 존속기간을 영구적인 것으로 약정하는 것도 허용된다.
  \item 지상권은 저당권의 목적이 될 수 있다.
  \item 지상권자는 지상권설정자의 명시적 반대에도 불구하고 지상권의 존속기간 내에서 그 토지를 타인에게 임대할 수 있다.
\end{choices}
\end{question}
\answer{2}
